% Generated draft: requires mathematical and source review.
\documentclass[11pt]{article}
\usepackage[margin=25mm]{geometry}
\usepackage{fontspec}
\setmainfont{Latin Modern Roman}
\usepackage{amsmath,amssymb,mathtools,unicode-math}
\setmathfont{Latin Modern Math}
\usepackage{enumitem,needspace,xurl,hyperref}
\setlength{\parindent}{0pt}
\setlength{\parskip}{5pt}
\setlength{\emergencystretch}{3em}
\newcommand{\N}{\mathbb{N}}
\newcommand{\Z}{\mathbb{Z}}
\newcommand{\Q}{\mathbb{Q}}
\newcommand{\R}{\mathbb{R}}
\newcommand{\C}{\mathbb{C}}
\newcommand{\F}{\mathbb{F}}
\newcommand{\A}{\mathbb{A}}
\newcommand{\PP}{\mathbb{P}}
\newcommand{\E}{\mathbb{E}}
\newcommand{\eps}{\varepsilon}
\newcommand{\dd}{\,\mathrm{d}}
\begin{document}
\section*{UnsolvedMath ID: 20001955}
\textbf{Source problem code:} catalogue entry
\subsection{Definitions and mathematical statement}
\textbf{Question 4.3, M.~Pflaum.} Is there a compact nonabelian Lie-group analogue of the Goldin--Holm--Knutson theorem? The supplied verification field identifies the preceding Theorem 4.1 as a ring-surjective map from inertial equivariant cohomology of a compact Hamiltonian torus space to Chen--Ruan cohomology of its regular symplectic quotient. Thus the original question asks for a suitable source product and ring theorem, not merely ordinary additive Kirwan surjectivity.

\textbf{Setting and precise additive partial statement from the aid.} Let a compact Lie group $G$, possibly disconnected, act in a Hamiltonian fashion on a compact symplectic manifold $(M,\omega)$, with equivariant moment map $\mu:M\to\mathfrak g^*$. Assume $0$ is a regular value, $Z=\mu^{-1}(0)\ne\varnothing$, and $X=[Z/G]$ is the resulting locally free quotient orbifold, retaining its finite stabilizer data. Put
\[
 M^g=\{x:gx=x\},\quad Z^g=Z\cap M^g,\quad
 C_G(g)=\{h:hg=gh\}.
\]
Let $\mathcal C_Z$ be the set of $G$-conjugacy classes $[g]$ with $Z^g\ne\varnothing$. Compactness and local freeness make this set finite; its elements have finite-order representatives. The corresponding inertia sectors are $[Z^g/C_G(g)]$. With rational coefficients, the aid asserts that restriction induces an ordinary-degree surjection
\[
 \bigoplus_{[g]\in\mathcal C_Z}
 H^*_{C_G(g)}(M^g;\Q)
 \longrightarrow
 \bigoplus_{[g]\in\mathcal C_Z}
 H^*_{C_G(g)}(Z^g;\Q).
\]
The target, after imposing the reduced-sector age shifts, is the additive vector space underlying $H^*_{\mathrm{CR}}(X;\Q)$. For each connected inertia sector $F$, its Chen--Ruan degree is ordinary degree plus $2\operatorname{age}_X(F)$, with age defined by stabilizer eigenvalues on the reduced tangent orbibundle. These shifts are not silently replaced by ages on $M$, and an ordinary-degree assertion is not presented as a canonical source-graded ring assertion.

The aid explicitly includes disconnected centralizers: the Kirwan map for $C_G(g)^0$ is claimed equivariant for the finite component group $\pi_0C_G(g)$, and taking its invariants is exact over $\Q$. This is the stated mechanism for extending the additive theorem to all compact $G$. It supplies no multiplicative compatibility with the standard Chen--Ruan obstruction-bundle product. The original ring-level nonabelian problem is therefore left separate.

\textbf{Source proposals and qualifications.} Pflaum motivates the question by quantization of singular reduced spaces and crossed-product algebras in symplectic orbifold theory. Goldin and Holm identify the source product as the main difficulty because the torus construction uses commutativity. Goldin, Holm and Pflaum expect an additive analogue and propose either indexing by conjugacy classes or replacing a fixed-locus factor by
\[
 H_G^*(U_g),\qquad U_g=\bigcup_{h\in G}M^{hgh^{-1}},
\]
in a product indexed by $g\in G$. These are proposed alternatives, not identified with the finite-sector direct sum above. Indeed the conjugacy-incidence space $G\times_{C_G(g)}M^g$ remembers which conjugate fixes a point, while its image $U_g\subset M$ can forget this information. No equivalence of their equivariant cohomologies is assumed.

The same contributors propose translating the torus result into proper \`etale Lie-groupoid language using [MM03], with a view toward a nonabelian product. Boyer warns about Morita-equivalent presentations and suggests Hochschild/cyclic homology of smooth crossed products or crepant resolutions as other approaches. The extracted expression $HC^*(C^\infty\mathbin{\rtimes}G)$ does not specify its underlying space, completion or cyclic theory, and is not promoted to a general isomorphism of Chen--Ruan rings. Nor are unprovided hypotheses for a crepant-resolution correspondence invented. The final Section 5 heading begins unrelated topological material and contains no further problem here.
\subsection{Short English statement}
Pflaum asks for a nonabelian extension of the ring-surjective orbifold Kirwan theorem. The supplied aid gives an additive rational surjection using finitely many conjugacy sectors and their centralizers, including disconnected groups. Defining a compatible source product, matching age gradings and obtaining a Chen--Ruan ring map remain separate requirements.
\subsection{Sources}
\begin{itemize}
\item Local UnsolvedMath registry, permanent ID 20001955, source code catalogue entry; canonical \texttt{aim-geometry-notes.json}, notes[292].
\item AIM \emph{Moment maps and surjectivity in various geometries}, Question 4.3 and Comments 4.4--4.7; supplied verification describes preceding Goldin--Holm--Knutson Theorem 4.1: \url{https://aimath.org/WWN/momentmaps/momentmaps.pdf}.
\item Question and quantization motivation: M.~Pflaum. Product issue: R.~Goldin and T.~Holm. Additive and groupoid proposals: R.~Goldin, T.~Holm and M.~Pflaum. Morita, cyclic-homology and crepant-resolution comments: C.~Boyer. Supplied label [MM03] refers to Moerdijk's groupoid approach; a full entry is absent.
\item Registry research attempt 1, dated 2026-08-14, partial result, rigor ``proved'', novelty confidence low. This explicitly machine-generated, non-peer-reviewed aid requires independent verification before citation; it supplies no separate proving citation for the additive extension.
\end{itemize}
\end{document}
